% Body text for the 4-page IEEE BigData Cup report.  Paste into the organisers'
% Overleaf template (kaggle_comp_report_fi_segm_2026.pdf) below its title block.
% Requires: amsmath, booktabs, graphicx, tikz.
%
% Every number here is produced by code in this repository and is reproducible
% from it.  Sections marked RESULTS-PENDING are filled from tools/compare_pods.py
% once the running ablations report; nothing else depends on them.

\section{Problem and approach}

We segment individual solar filaments in 2048$\times$2048 GONG H$\alpha$
full-disk images. The training set is 707 observations carrying 1154 annotator
readings; 42\% of observations are read by more than one annotator, and those
annotators agree only moderately. The evaluation set is 180 images.

Our system is a U-Net decoder over an ImageNet-pretrained EfficientNet-B0
encoder, trained on 512-pixel tiles and applied to whole observations by
cosine-blended tiling. Three choices are specific to this problem rather than
inherited from a standard recipe.

\textbf{Geometry-aware preprocessing.} The solar disk is fitted per image and
limb darkening removed by a radial flat field, so that a filament of given
contrast presents the same way at disk centre and near the limb. Features are
standardised against constants fixed once across the corpus, and every
off-disk pixel is excluded from both loss and metric.

\textbf{A decoder that reaches full resolution.} Standard timm encoders stride
the input by 32 and the usual U-Net stops decoding at stride 2, upsampling
logits at the end. Filament barbs are a few pixels wide, so our decoder runs
back to stride 1 with learned blocks and deep supervision at strides 2 and 4.

\textbf{A loss matched to a stochastic target.} Because a third of images carry
multiple readings, the per-pixel target is genuinely random: the correct thing
to predict is $P(\text{a randomly drawn annotator marks this pixel})$. Binary
cross-entropy is a proper scoring rule, so its minimiser \emph{is} that
posterior, which is what makes the later threshold search meaningful. Soft Dice
is added at low weight purely as a conditioner on the 0.3\% positive rate.

Instances are extracted by hysteresis thresholding (a strict seed level, a
looser extent level) followed by connected components, with morphological
closing and area filters. Every one of those decision parameters is fitted
against Panoptic Quality itself rather than against a proxy.

\section{Validation methodology}

The competition metric is Panoptic Quality,
\begin{equation}
\mathrm{PQ} = \frac{\sum_{(p,g)\in \mathrm{TP}} \mathrm{IoU}(p,g)}
                   {|\mathrm{TP}| + \tfrac12|\mathrm{FP}| + \tfrac12|\mathrm{FN}|},
\end{equation}
matching at $\mathrm{IoU} > 0.5$ strictly and pooled over every
annotator--image entry. We re-implemented the organisers' rules independently
and assert agreement between the two implementations on hand-computed cases,
multi-annotator observations, images with no ground truth and images with no
prediction (7 automated checks). We did this because two of our own early
conclusions turned on metric details --- fragmentation counting uses
\emph{any} overlap, not a threshold --- and a metric taken on trust is a
result taken on trust.

Splits are five-fold and stratified by observing site, so that a fold never
scores an instrument it has also trained on. Two further precautions matter:

\begin{itemize}
\item \textbf{Splits separate observations, not readings.} An observation read
by three annotators yields three scoring entries. Splitting on entries put
59.8\% of nominally held-out observations into the tuning half as well. This
was a real defect in our own earlier work, found by a test that now guards it.
\item \textbf{Operating points are fitted and reported on disjoint halves.}
Fitting the decision parameters on the same data they are reported on inflates
PQ; measured directly, that selection optimism is $+0.0024$ for our pooled fit.
\end{itemize}

\section{Results}

Five-fold cross-validation with per-fold operating points gives
$\mathrm{PQ} = 0.3946 \pm 0.0113$. Pooling the out-of-fold predictions and
fitting one operating point over all of them gives \textbf{PQ 0.4167},
reported on 353 observations disjoint from the 354 used to fit. Our public
leaderboard score is 0.37.

\begin{table}[t]
\centering
\caption{Out-of-fold error decomposition, 707 observations.}
\begin{tabular}{lr@{\hskip 2em}lr}
\toprule
PQ & 0.4167 & one-to-many & 144 \\
SQ (mean matched IoU) & 0.6748 & many-to-one & 70 \\
RQ & 0.6175 & mean matched Dice & 0.8027 \\
true positives & 2501 & mean semantic Dice & 0.6517 \\
false positives & 1458 & precision & 0.632 \\
false negatives & 1640 & hit rate & 0.604 \\
\bottomrule
\end{tabular}
\end{table}

\subsection{What limits the score}

We measured the achievable ceiling directly, by scoring each annotator's own
masks against every reading of the same image under the real metric. A method
emitting human-quality masks would score \textbf{0.7687} overall: 1.0 on the
411 single-annotator observations, 0.6307 on the 296 multi-annotator ones.

It is tempting to conclude that annotator disagreement is what bounds us. It is
not, and the evidence is direct: out-of-fold PQ on single-annotator
observations, where an exact match is achievable by construction, is
\textbf{0.4163}, against \textbf{0.4169} on multi-annotator observations. If
disagreement bound us we would score far better where there is none.

The binding constraint is mask quality, and it is remarkably rigid. Mean
matched IoU reaches $\approx 0.67$ by epoch six and does not move for the
remaining 24 epochs, while PQ swings by 0.13 through recognition alone. Every
run we have trained --- both encoders, all five folds, 20-epoch and 30-epoch
schedules --- lands between 0.63 and 0.68. Exactly one matched pair in 2501
exceeds IoU 0.90.

That rigidity matters because the metric is unusually sensitive to it. Shifting
every instance's IoU by a constant and recomputing PQ gives a near-linear
response of \textbf{1.06 PQ per unit of mean IoU} (Table~\ref{tab:headroom}) ---
so $+0.05$ IoU is worth more than converting \emph{every} one of our 321
near-miss instances (unmatched ground truth in IoU 0.4--0.5), which is worth
$+0.048$. Predictions are only 8\% too large by area, far too little to explain
an IoU of 0.67, so the error is in \emph{where} the boundary sits rather than
how large the mask is.

\begin{table}[t]
\centering
\caption{PQ under a uniform shift in every instance's IoU.}
\label{tab:headroom}
\begin{tabular}{rrr}
\toprule
IoU shift & PQ & $\Delta$ \\
\midrule
$-0.05$ & 0.3597 & $-0.057$ \\
$+0.02$ & 0.4391 & $+0.022$ \\
$+0.05$ & 0.4700 & $+0.053$ \\
$+0.10$ & 0.5225 & $+0.106$ \\
\bottomrule
\end{tabular}
\end{table}

\subsection{Ablations}

% RESULTS-PENDING: filled from tools/compare_pods.py.  Each row is measured
% against its own baseline trained on the same card, in the same precision,
% under the same code version and seed; the noise floor is the spread between
% identical baseline runs and is quoted so the reader can judge significance.

\section{Negative results}

We report these because they were expensive to establish and because two of
them corrected conclusions we had previously drawn.

\textbf{Label smoothing (0.05) costs $-0.011$.} Against an older code version it
appeared to gain $+0.022$. It does not; the apparent gain was a comparison
across code versions.

\textbf{Ensemble threshold calibration is worth nothing measurable.} Our
operating point is fitted on single-model out-of-fold maps but applied to a
five-model average, and averaging shaves probability peaks, so the same
threshold admits less. We built a label-free correction that holds admitted
pixel mass fixed rather than the threshold, and measured it on the test images'
probability histograms (test pixels, never test labels). It moves
$\text{seed} = 0.950 \rightarrow 0.929$, raises predicted instances per image
from 6.08 to 6.43 as intended, and scored \emph{identically} to the
uncalibrated submission. The operating point was not the problem.

\textbf{Capacity, as previously measured, was confounded.} EfficientNet-B4
scored $-0.021$ against B0 --- but on a Turing T4 in fp16, where B4 overflowed
and the gradient scaler silently dropped the offending steps. That run cannot
distinguish ``worse'' from ``undertrained''. We re-tested it on Ampere hardware
in bf16, which carries fp32's exponent range and cannot overflow that way.

\section{Strengths, weaknesses, and what we would do next}

\textbf{Strengths.} The metric is independently verified rather than assumed.
Every reported number comes from a split that separates observations, with
selection optimism measured rather than hoped away. The achievable ceiling is
measured, not argued. Ablations are paired within a single machine, precision
and code version, because a difference measured across any of those is not a
difference.

\textbf{Weaknesses.} Mean matched IoU is the binding constraint and we have not
moved it. Fragmentation persists (144 one-to-many, 70 many-to-one). Ten of 707
observations emit nothing at all, which is certainly wrong: no training reading
contains zero filaments. And a $-0.047$ gap between cross-validation and the
leaderboard remains unexplained now that calibration is excluded.

\textbf{Next.} The objective does not optimise what the metric measures: our
Dice term is semantic, pooled over the whole tile, while PQ scores per-instance
IoU. A boundary-weighted term, applied as a per-pixel re-weighting of BCE so
that the proper-scoring-rule property survives, is the direct test and costs
nothing to reject. Beyond that, the auxiliary spine head --- explicitly
sanctioned by the organisers, and targeting the structural continuity they name
as a core challenge --- and self-supervised pretraining on unlabelled GONG
H$\alpha$ are the two levers with the most room left.

\section{Reproducibility}

All code, tests and tooling are at \url{https://github.com/msraaghavan/Solar}
under MIT, with \texttt{requirements.txt} pinning the runtime the submitted
results were produced on. The pipeline notebook is generated from the same
source it documents, so it cannot drift from it. The test suite is 51 checks
and needs no GPU; several of them encode defects found in this project's own
earlier work, including the reading-level split leak, a censored tuning grid,
and a fragmentation rule that used a threshold where the organisers' uses any
overlap.
